\subsection{Применяемые технологии}

Разработка программного средства обнаружения аномалий в Kubernetes-кластере осуществлялась с использованием современного стека технологий, ориентированного на создание масштабируемых распределённых систем и обработку данных с применением методов машинного обучения.

В качестве основного языка программирования был выбран \texttt{Python} \cite{Python}, который широко применяется при разработке систем анализа данных и машинного обучения благодаря развитой экосистеме библиотек и инструментов.

Для реализации серверной части приложения использован фреймворк \texttt{FastAPI} \cite{FastAPI}, обеспечивающий высокую производительность за счёт поддержки асинхронного ввода-вывода и удобных инструментов для разработки REST API. Запуск приложения осуществляется с помощью ASGI-сервера \texttt{Uvicorn} \cite{Uvicorn}, позволяющего эффективно обрабатывать большое количество одновременных запросов.

Для работы с базой данных применена библиотека \texttt{SQLAlchemy} \cite{SQLAlchemy}, реализующая объектно-реляционное отображение и обеспечивающая удобное взаимодействие с реляционной СУБД \texttt{PostgreSQL} \cite{PostgreSQL}. В качестве драйвера подключения используется \texttt{asyncpg} \cite{asyncpg}, что позволяет организовать асинхронную работу с базой данных и повысить общую производительность системы.

Организация фоновых вычислений реализована с использованием системы распределённых задач \texttt{Celery} \cite{Celery}. Данный инструмент позволяет выносить ресурсоёмкие операции, такие как обучение нейросетевых моделей и выполнение задач детектирования, в отдельные worker-процессы. Это обеспечивает асинхронность выполнения и возможность горизонтального масштабирования системы.

Для реализации модели обнаружения аномалий использована библиотека \texttt{PyTorch} \cite{PyTorch}, предоставляющая гибкие средства для построения и обучения нейронных сетей, включая рекуррентные архитектуры, такие как LSTM. Для выполнения вспомогательных задач предобработки данных и оценки качества моделей применяется библиотека \texttt{scikit-learn}, широко используемая для решения задач машинного обучения \cite{ScikitLearn}.

В качестве источника данных используется система \texttt{OpenSearch} \cite{OpenSearch}, из которой извлекаются audit-логи Kubernetes-кластера. Для взаимодействия с данным хранилищем применяется клиентская библиотека \texttt{opensearch-py} \cite{OpenSearchPy}, обеспечивающая выполнение запросов и получение данных для последующего анализа.

Выбранный стек технологий обеспечивает необходимый уровень производительности, расширяемости и удобства разработки, а также позволяет эффективно реализовать распределённую систему обнаружения аномалий, работающую с потоками событий в Kubernetes-кластере.